\documentclass[12pt, a4paper]{article}

\usepackage{graphicx}

\textheight=210mm
\textwidth=155mm

\hoffset=-1mm
\voffset=-10mm

%% for corrections 
\usepackage{color}
\newcommand\bb{\begin{color}{blue}}
\newcommand\br{\begin{color}{red}}
\newcommand\eb{\end{color}\ }

\begin{document}

\begin{center}
{\bf Response to Referee~\#1's comments}\\
on paper ``Power spectrum scaling as a measure of critical slowing down''\\[5pt]
by {\it J.~Prettyman}
\end{center}

\vspace{3mm}

\noindent We would like to thank the Referee for their time in reviewing this manuscript and for the comments which have informed the revisions made. 

We agree that the paper, as originally presented, is mostly theoretical and lacking in specific content to place it within the context of Earth-system tipping-point research. Our own previous research [Prettyman et al. 2019]\footnote{Prettyman, J., et al., 2019. \emph{Generalized early warning signals in multivariate and gridded data with an application to tropical cyclones.} Chaos. 073105.} focussed on applying this method to identifying Critical Slowing Down in sea-level pressure time series data close to hurricane landfall sites. It is as an extension to this research that we present this new paper: a justification of the method by considering the analytical properties of the critical slowing down phenomenon. We hope, in future, to apply this method to more systems, particularly coral bleaching, which we mention in the revised manuscript (page 6). It is also our hope that by publishing this method here, this method will become a useful tool in the tipping-point analysis canon, and that more applications may be found. 

Whilst the revision time for this manuscript did not allow for new research, we have included an example application (page 8) to paleo-temperature proxy records which have previously been studied using other methods [Lenton et al. 2012]\footnote{Lenton, T.M. et al., 2012. \emph{Climate bifurcation during the last deglaciation?}. Climate of the Past, 8(4), pp.1127-1139.}. The PS indicator confirms the presence of critical slowing down in the time series but is not able to provide an early warning signal, which motivates future research in areas where high time-resolution datasets are available (such as modern meteorological datasets). 

In accordance with the comments made by the referees, the manuscript has been updated with a view to placing it within the context of earth-system tipping point research.

\vspace{7mm}

\noindent
We hope that the revised manuscript is now suitable for publication
in the ERL special issue.

\vspace{5mm}

\noindent
Yours sincerely,

\vspace{1mm}

\noindent
T.~Kuna, V.~N.~Livina, and J.~Prettyman

\end{document}

